\section{Introduction}

Fine-tuning is the dominant method for adapting large language models to
downstream tasks and deployment contexts.
Yet the internal computational structures that form during pre-training---the
\emph{circuits} that implement specific algorithmic capabilities---are rarely
monitored during fine-tuning.
If safety-relevant circuits degrade or reorganise through this process, the
resulting model may pass capability benchmarks while having silently lost
properties present in the base model.

Induction heads provide a concrete, well-characterised target for this
investigation.
\citet{olsson2022context} showed that a two-head circuit---a previous-token
head in layer~0 that copies token identities into the residual stream, followed
by an induction head in layer~1 that reads this information to attend back to
previous occurrences---is responsible for the majority of in-context learning
in small transformers.
This circuit is mechanistically understood: its causal role has been verified
via activation patching~\citep{conmy2023automated}, and analogous circuits are
present across architectures and scales~\citep{elhage2021mathematical}.

\textbf{Our question:}
Does this circuit survive fine-tuning on a narrow distribution, and if it
degrades or reorganises, at what training step does the transition occur?

\textbf{Contributions:}
\begin{itemize}[leftmargin=*]
  \item We replicate the induction circuit in \texttt{attn-only-2l} and verify
    its causal role via the activation-patching protocol of
    \citet{conmy2023automated}.
  \item We fine-tune on 500K tokens of Python code (primary) and TinyStories
    prose (mandatory control), saving checkpoints every 100 steps and computing
    all eight circuit metrics from Section~\ref{sec:methods} at each checkpoint.
  \item \textbf{Preliminary finding (seed 42; 3 checkpoints).} We find that
    the L1H6 induction circuit's prefix-matching score \emph{increases} under
    both code and prose fine-tuning ($0.408 \to 0.602 \to 0.650$ for code;
    $0.408 \to 0.583 \to 0.617$ for prose, over steps $0/100/200$)---a result
    not reported in any prior work---which implies that fine-tuning on a
    narrow distribution does not necessarily erode the induction circuit, and
    that safety evaluations for fine-tuned models should include circuit-level
    diagnostics to detect such changes (whether degradation or reinforcement)
    rather than assuming degradation by default.
  \item We release a fully reproducible codebase with interactive dashboard
    (Hugging Face Spaces) and all experiment configs.
\end{itemize}

The rest of this paper is organised as follows.
Section~\ref{sec:background} reviews the transformer circuits framework and
activation-patching methodology.
Section~\ref{sec:methods} specifies the experimental protocol.
Section~\ref{sec:results} presents the main findings.
Section~\ref{sec:safety} discusses safety implications.
Section~\ref{sec:limitations} states limitations.
Section~\ref{sec:conclusion} concludes.
